\section{API Specification}

\subsection{Contract Strategy}
AgriApp adopts \texttt{/api/v1} as the canonical API namespace for all active client integrations (web and Android). The objective is contract stability across iterative releases: new features are added under versioned routes, while historical aliases are retained only for backward compatibility during migration.

All operational clients should target \texttt{/api/v1} endpoints first. Legacy aliases have been removed from the active runtime surface.

\subsection{Canonical API v1 Endpoints}
\begin{longtable}{p{0.36\linewidth} p{0.14\linewidth} p{0.45\linewidth}}
\toprule
Path & Method & Purpose \\
\midrule
\texttt{/api/v1/health} & GET & Service liveness and API version marker \\
\texttt{/api/v1/health/checklist} & GET & Live checklist (RPi, camera, GPS, BME, storage; ESP backend field kept available) \\
\texttt{/api/v1/system/status} & GET & Hostname, storage info, runtime timestamp \\
\texttt{/api/v1/system/server/restart} & POST & Restart \texttt{agriapp-server.service} (user mode) \\
\texttt{/api/v1/system/reboot} & POST & Reboot Raspberry Pi (requires sudo policy) \\
\texttt{/api/v1/system/poweroff} & POST & Power off Raspberry Pi (requires sudo policy) \\
\texttt{/api/v1/network/mode} & GET/POST & Read/set network mode (\texttt{wifi\_only}, \texttt{ap\_only}, \texttt{hybrid\_debug}) \\
\texttt{/api/v1/logs} & GET & Unified pull-based runtime logs (server + loops) \\
\texttt{/api/v1/logs/stream} & GET & Unified push log stream (SSE) \\
\texttt{/api/v1/images} & GET & Paginated image listing + metadata/filtering \\
\texttt{/api/v1/images/\{filename\}/status} & GET & Per-image status, file size, metadata summary \\
\texttt{/api/v1/images/\{filename\}/thumbnail} & GET & Cached thumbnail stream (query \texttt{w}) \\
\texttt{/api/v1/images/delete} & POST & Delete one image and related artifacts \\
\texttt{/api/v1/images/delete-all} & POST & Delete all images and related artifacts \\
\texttt{/api/v1/images/upload} & POST & Upload image from external producer \\
\texttt{/api/v1/labels} & GET & Read label state for one image \\
\texttt{/api/v1/labels} & POST & Save labels (JSON + YOLO export sync) \\
\texttt{/api/v1/capture/rpi/oneshot} & POST & Trigger one Raspberry Pi camera capture \\
\texttt{/api/v1/capture/esp/oneshot} & GET & Trigger ESP capture via backend proxy \\
\texttt{/api/v1/esp/status} & GET & Probe ESP status endpoint via backend \\
\texttt{/api/v1/capture/loop/status} & GET & Inspect unattended capture loop state \\
\texttt{/api/v1/capture/loop/start} & POST & Start unattended capture loop (\texttt{interval\_seconds}) \\
\texttt{/api/v1/capture/loop/stop} & POST & Stop unattended capture loop \\
\texttt{/api/v1/capture/esp/loop/status} & GET & Inspect unattended ESP capture loop state \\
\texttt{/api/v1/capture/esp/loop/start} & POST & Start unattended ESP loop (\texttt{interval\_seconds}, \texttt{esp\_base}) \\
\texttt{/api/v1/capture/esp/loop/stop} & POST & Stop unattended ESP capture loop \\
\texttt{/api/v1/download/dataset} & GET & Download full dataset archive \\
\texttt{/api/v1/download/dataset-selected} & POST & Download archive for selected filenames \\
\texttt{/api/v1/images/download-with-labels} & GET & Download one image bundle (image+labels+json+metadata) \\
\bottomrule
\end{longtable}

\subsection{UI Routes}
\begin{longtable}{p{0.36\linewidth} p{0.14\linewidth} p{0.45\linewidth}}
\toprule
Path & Method & Purpose \\
\midrule
\texttt{/} & GET & Entry page \\
\texttt{/gallery} & GET & Web gallery interface \\
\texttt{/capture} & GET & Web capture/control page \\
\texttt{/system} & GET & Web system/network page \\
\texttt{/health} & GET & Web live checklist page (RPi/camera/GPS/BME/storage) \\
\texttt{/log} & GET & Web runtime log page \\
\texttt{/label} & GET & Web labeler interface (query \texttt{image}) \\
\bottomrule
\end{longtable}

\subsection{Compatibility (Legacy) Routes}
Legacy aliases have been removed from the runtime surface after migration of active clients to \texttt{/api/v1}. New integrations must target canonical versioned endpoints only.

\subsection{Current Web Exposure Policy}
The web frontend currently runs in RPi-only control mode for field stability. This means ESP controls are intentionally hidden in web routes (\texttt{/capture}, \texttt{/system}, \texttt{/health}), while ESP APIs remain available in \texttt{/api/v1} for diagnostics and phased reintegration.

\subsection{Endpoint Grouping by Device}
\subsubsection{Raspberry Pi Edge Node}
Core operational routes for the edge runtime include system status/control, network mode switching, local capture, ESP proxy capture, and unattended loop control. These endpoints are consumed by both browser and tablet clients, and are also used for shell-based validation via \texttt{curl}.

\subsubsection{Web Client}
The web interface consumes the same \texttt{/api/v1/images}, \texttt{/api/v1/labels}, deletion, thumbnail, and dataset download contracts used by other clients. This avoids browser-specific backend branches and keeps semantics aligned with tablet behavior.

\subsubsection{Android Tablet Client}
The Android app focuses on operational routes: discovery + health, status, capture one-shot (RPi/ESP), capture loop start/stop (RPi and ESP), gallery pagination with thumbnails, and system/network actions. The app intentionally does not introduce private backend APIs; it reuses the same versioned surface.

\subsection{Request and Response Semantics}
Most endpoints return JSON payloads with explicit \texttt{status} semantics. Binary endpoints (image and ZIP downloads, thumbnail streams, ESP capture proxy stream) return bytes and may include additional headers.

Pagination is server-side and stable: responses include \texttt{page}, \texttt{page\_size}, \texttt{total\_items}, and \texttt{total\_pages}. This enables consistent navigation rules between web and tablet UIs.

Image items include contextual metadata fields (GPS latitude/longitude and BME280-derived temperature/humidity/pressure when available). Missing sensor readings are represented as \texttt{null} and should be rendered as \texttt{N/A} in client UIs.

\subsection{Capture-Specific Notes}
\begin{itemize}
  \item \textbf{RPi one-shot} (\texttt{/api/v1/capture/rpi/oneshot}) executes \texttt{capture\_rpi.sh --oneshot}, stores the generated image, enriches metadata, and returns latest filename.
  \item \textbf{ESP one-shot} (\texttt{/api/v1/capture/esp/oneshot}) requires query \texttt{esp\_base}, validates private IP constraints, executes \texttt{capture\_esp.py}, stores an \texttt{esp\_...} image locally, and streams JPEG bytes as response. The response includes optional headers carrying saved filename and metadata hints (\texttt{X-AgriApp-*}) for client-side fast feedback.
  \item \textbf{RPi loop} endpoints (\texttt{/api/v1/capture/loop/*}) control unattended periodic capture via backend-managed process state; \texttt{/start} accepts \texttt{interval\_seconds} in a bounded range.
  \item \textbf{ESP loop} endpoints (\texttt{/api/v1/capture/esp/loop/*}) expose the same control model for ESP, with mandatory \texttt{esp\_base} on start. The loop invokes the backend ESP one-shot proxy path to preserve metadata sidecar consistency across manual and unattended captures.
\end{itemize}

\subsection{Error Model}
The API uses conventional HTTP statuses:
\begin{itemize}
  \item \texttt{2xx}: success,
  \item \texttt{4xx}: request or state errors (missing params, invalid mode, resource not found),
  \item \texttt{5xx}: backend/runtime execution failures.
\end{itemize}

In failure paths, JSON responses include compact diagnostic context (e.g., command stderr tail, validation message, timeout cause) to support field troubleshooting without requiring immediate shell access.

\subsection{Migration Guidance}
Client integrations are now converged on \texttt{/api/v1}. Future API evolution should preserve versioned routing discipline (\texttt{/api/vN}) and avoid introducing unversioned operational aliases.
